\documentclass[10pt,a4paper]{amsart}
\usepackage[margin=19mm]{geometry}
\usepackage{amsmath,amssymb,mathtools}
\usepackage[scr=rsfs,cal=euler]{mathalfa}
\usepackage{xfrac}
\usepackage[unicode,hidelinks,bookmarks=false]{hyperref}
\newcommand{\ie}{i.e.\ }
\newcommand{\cf}{cf.\ }
\newcommand{\resp}{respectively\ }
\newcommand{\Subsection}{\subsection}
\providecommand{\nobreakdash}{\nobreak\hbox{-}}
\setlength{\emergencystretch}{2em}
\multlinegap0pt
\usepackage{xcolor}
\usepackage{amssymb}
\usepackage{float}
\usepackage{tikz}
\usepackage[shortlabels]{enumitem}
\newtheorem{lemma}{Lemma}
\newtheorem{theorem}{Theorem}
\usepackage{cleveref}
\crefname{lemma}{Lemma}{Lemmas}
\crefname{theorem}{Theorem}{Theorems}
\DeclareMathOperator{\LS}{RP}
\DeclareMathOperator{\ROS}{ROS}
\title{Further results on latin squares with disjoint subsquares using rational outline squares}
\author{Tara Kemp, James Lefevre}
\date{Source: arXiv:2505.07252v1 (CC BY 4.0)}
\begin{document}
\maketitle

\noindent\emph{Source and licence.} Further results on latin squares with disjoint subsquares using rational outline squares. Version arXiv:2505.07252v1 (CC BY 4.0), distributed by arXiv under the Creative Commons Attribution 4.0 International licence, http://creativecommons.org/licenses/by/4.0/, as recorded in the arXiv licence metadata for this exact version and shown on the abstract page https://arxiv.org/abs/2505.07252v1. Full licence notice: \texttt{licenses/arxiv-2505.07252-CC-BY-4.0.txt}.\par\smallskip

\noindent\textit{Editorial scope.} Counted unit: the article's theorem sufiĉa? (source0.tex lines 408--413). The article's complete original proof of it is delivered with it (source0.tex lines 414--439, one \texttt{proof} environment of 26 lines). No mathematics is written, repaired or re-derived: the statement and the proof are the article's own lines, in the article's own notation, and no retained line is substituted or re-flowed.  Retained uncounted as the source-local context the proof needs: lines 355--358.  Nothing was omitted from any retained span, so the package carries no omission record.  No retained line contains a citation, so the wrapper carries no bibliography and no bibliography entry.  No retained line contains a figure environment, an image-inclusion command or drawing code, so no image is delivered and the package declares no asset dependency.  Editorial formatting only, in addition: an AMS \texttt{amsart} preamble that restates the article's own declarations and macros used by the retained lines, namely the environments \texttt{lemma}, \texttt{theorem} on separate counters, together with the standard packages those lines need.  The retained environments print in this document as Lemma 1, Theorem 1 in order of appearance, whereas the article prints the same items with the numbering of its own full document.  The complete native source of the article as distributed in the arXiv e-print for this exact version is bundled unchanged as \texttt{sources/177946/source0.tex}.
\begin{lemma}
\label{lemma: always symmetric solution}
    If an $\LS(h_1\dots h_k)$ exists, then an $\ROS(h_1\dots h_k)$ exists.
\end{lemma}

\begin{theorem}
\label{condition: sufiĉa?}
    If an $\LS(h_1\dots h_k)$ exists, then
    $$\left(\sum_{i\in A\cup C}h_i\right)^2 + \left(\sum_{i\in B\cup D}h_i\right)^2 - \sum_{i\in E}h_i^2 \geq \left(\sum_{i\in A\cup D}h_i\right)\left(\sum_{i\in B\cup C}h_i - \sum_{j\in \overline{E}}h_j\right),$$
    where $A$, $B$, $C$ and $D$ are pairwise disjoint subsets of $[k]$, $E = A\cup B\cup C\cup D$ and $\overline{E} = [k]\setminus E$.
\end{theorem}

\begin{proof}
    By \Cref{lemma: always symmetric solution}, if such a latin square exists, then an $\ROS(h_1\dots h_k)$ exists also. We consider this outline square $X$.

    Observe that there are $(\sum_{i\in A\cup D}h_i)(\sum_{i\in B\cup C}h_i)$ symbols, including repetition, required across the cells $(i,j)$ where $i\in A\cup D$ and $j\in B\cup C$. Of these symbols, at most $(\sum_{i\in A\cup D}h_i)(\sum_{j\in \overline{E}}h_j)$ can be symbols in $\overline{E}$. The remaining symbols must be from $E$.

    We now count the number of symbols from $E$ in these cells. First, note that there are at most $(\sum_{i\in A}h_i)^2 - \sum_{i\in A}h_i^2$ symbols from $A$ in the rows of $A$, and likewise for the rows of $D$ and columns of $B$ and $C$. Thus, we must now consider the symbols appearing not within columns or rows from the same subset. The number of symbols from $A$ within rows of $D$ of these cells is given by
    $$\sum_{a\in A} \sum_{d\in D} \sum_{i\in B\cup C} X(a,d,i)$$
    and the number of symbols from $D$ in columns of $A$ is the same by symmetry. Similarly, the number of symbols from $B$ in columns of $C$ and the number of symbols from $C$ in columns of $B$ is given by
    $$\sum_{b\in B} \sum_{c\in C} \sum_{i\in A\cup D} X(b,c,i).$$

    Combining these, and using symmetry to change the order of the terms in the $X$ variables, we see that
    \begin{align*}
        2\sum_{a\in A} \sum_{d\in D} \sum_{i\in B\cup C} X(a,d,i) + 2\sum_{b\in B} \sum_{c\in C} \sum_{i\in A\cup D} X(b,c,i) &= 2 \sum_{a\in A} \sum_{c\in C} \sum_{i\in B\cup D} X(a,c,i) + 2 \sum_{b\in B} \sum_{d\in D} \sum_{i\in A\cup C} X(b,d,i)\\
        &\leq 2 \sum_{a\in A} \sum_{c\in C} \sum_{i\in E} X(a,c,i) + 2 \sum_{b\in B} \sum_{d\in D} \sum_{i\in E} X(b,d,i).
    \end{align*}
    The two terms in this last expression are the number of symbols in cells $(i,j)$ where $i\in A$ and $j\in C$ or $i\in B$ and $j\in D$ respectively. Thus, we have that
    $$2 \sum_{a\in A} \sum_{d\in D} \sum_{i\in B\cup C} X(a,d,i) + 2 \sum_{b\in B} \sum_{c\in C} \sum_{i\in A\cup D} X(b,c,i) \leq 2(\sum_{i\in A}h_i)(\sum_{i\in C}h_i) + 2(\sum_{i\in B}h_i)(\sum_{i\in D}h_i).$$

    Altogether, the number of possible entries is at least the number of required symbols, so
    \begin{align*}
    \begin{split}
        \left(\sum_{i\in A\cup D}h_i\right)\left(\sum_{i\in B\cup C}h_i - \sum_{j\in \overline{E}}h_j\right)\leq \left(\sum_{i\in A}h_i\right)^2 + \left(\sum_{i\in B}h_i\right)^2 + \left(\sum_{i\in C}h_i\right)^2 + \left(\sum_{i\in D}h_i\right)^2 - \sum_{i\in E}h_i^2\\ \phantom{1} + 2\left(\sum_{i\in A}h_i\right)\left(\sum_{i\in C}h_i\right) + 2\left(\sum_{i\in B}h_i\right)\left(\sum_{i\in D}h_i\right)
    \end{split}
    \end{align*}
    $$\left(\sum_{i\in A\cup D}h_i\right)\left(\sum_{i\in B\cup C}h_i - \sum_{j\in \overline{E}}h_j\right)\leq \left(\sum_{i\in A\cup C}h_i\right)^2 + \left(\sum_{i\in B\cup D}h_i\right)^2 - \sum_{i\in E}h_i^2.$$
\end{proof}

\end{document}
